%% kapa: exam=1 round=21 year=2010 period=1 form=A subject=economics subject_ko=경제원론 seq=23 number=63
%% source: Q-Net 감정평가사 기출문제 (cst00309) artlSeq=5024181 file=extracted/1008/5024181/2010년 제21회 감정평가사 1차시험 기출문제/1교시.hwp
%% conversion: kordoc 4.24.51 -> tools/gen1.py (auto); review=proofread-sonnet; flags=tikz
\begin{question}{63}
다음 공급곡선들의 각 점에서 측정한 가격탄력성 크기를 옳게 비교한 것은?
\begin{center}
\begin{tikzpicture}[font=\small, x=1cm, y=1cm]
  \draw (0,0) -- (0,4.6);
  \draw (0,0) -- (4.8,0);
  \node[above] at (0.1,4.6) {\textbf{가격}};
  \node[below] at (4.7,-0.1) {\textbf{수량}};
  \draw (0,0) -- (1.65,3.63);
  \draw (0,0) -- (3.39,3.04);
  \draw (0,0) -- (4.21,1.41);
  \fill (1.03,2.27) circle (2pt);
  \fill (1.96,1.77) circle (2pt);
  \fill (2.67,0.90) circle (2pt);
  \node at (0.78,2.29) {\textit{A}};
  \node at (1.72,1.81) {\textit{B}};
  \node at (2.57,1.15) {\textit{C}};
\end{tikzpicture}
\end{center}
\begin{choices}
  \item A > B > C
  \item C > B > A
  \item A > B = C
  \item C = B > A
  \item A = B = C
\end{choices}
\end{question}
\answer{5}
